\documentclass[11pt]{article}

\usepackage[T1]{fontenc}
\usepackage[margin=1in]{geometry}
\usepackage{amsmath,amssymb,mathtools,amsthm}
\usepackage{booktabs,tabularx,array}
\IfFileExists{lmodern.sty}
  {\usepackage{lmodern}\usepackage{microtype}}
  {\usepackage[expansion=false]{microtype}}
\usepackage{enumitem}
\usepackage[hidelinks]{hyperref}

\setlength{\parindent}{0pt}
\setlength{\parskip}{0.55em}
\setlist[itemize]{leftmargin=*,itemsep=0.15em,topsep=0.25em}
\renewcommand{\arraystretch}{1.12}

\newcommand{\AP}{\operatorname{AP}}
\newcommand{\CNAP}{\operatorname{CNAP}}
\newcommand{\Regime}{\mathcal{R}}
\newcommand{\Assess}{\mathcal{A}}
\newcommand{\CSet}{\mathcal{C}}
\newcommand{\PSet}{\Pi_\star}
\newcommand{\E}{\mathbb{E}}
\newcommand{\Prb}{\mathbb{P}}
\newcommand{\APta}{\widetilde{\AP}}
\newcommand{\CNAPta}{\widetilde{\CNAP}}

\newtheorem{definition}{Definition}
\newtheorem{proposition}{Proposition}
\theoremstyle{remark}
\newtheorem*{remark}{Remark}

\title{Supported Model Replacement Across Target Prevalences\\
\large A Paired Replication-Based Method for Average Precision}
\author{Marcus A. Thomas}
\date{Methodological proposal}

\begin{document}

\maketitle

\begin{abstract}
A decision about whether to replace an incumbent model with a candidate often begins with an observed difference in a chosen performance measure. When the candidate's advantage is intended to support replacement across target conditions or empirical evaluations, a further question arises: how much of that advantage remains when the replacement claim is subjected to those declared challenges? We address this question with a falsification-oriented support construction defined on paired, condition-indexed evaluation effects.

The construction yields two complementary summaries: supported magnitude, which describes the effect retained under the challenge, and literal survival of the claimed advantage, which records its persistence above a declared floor. The condition set, evaluation regime, and challenge order determine the scope and severity of the assessment.

We develop the construction for model replacement using a prior-standardized and chance-normalized form of average precision. The framework separates evidence supplied by distinct evaluations from computational projections based on resampling a single evaluation. The proposed finite-evidence replacement rule requires supported magnitude and literal survival to exceed prespecified thresholds. The construction can be applied to other performance measures, such as AUROC, once their relevant effects and conditions of challenge have been specified.
\end{abstract}


\section{Performance Claims and their Support}
\label{sec:introduction}

Prediction systems often rank cases so that a selected set can receive further
attention. A radiology service, for example, may send its highest-scored images
for review. The composition of that set matters, not only the separation
between outcome classes, because the selected set is what the decision maker
receives.

Separation (eg AUROC) and selected-set purity (eg Average Precision) part company when prevalence changes.
Suppose a threshold selects $80\%$ of positive cases and $5\%$ of negative
cases. In a balanced population of 10,000 cases, it returns 4,000 positives and
250 negatives, for precision $0.94$. In a population with $1\%$ positives, the
same threshold returns 80 positives and 495 negatives, for precision $0.14$.
Its true-positive and false-positive rates have not moved. The supply of
negative cases has, and that is enough to change the meaning of the selected
set.

Average precision aggregates precision as the ranking is traversed, so ordinary
empirical AP evaluates performance using the positive fraction in the
evaluation data. Let $D$ denote those data: the observed outcomes and model
scores on a common set of $N$ evaluation units. Thus $N$ is the number of data
points on which AP is calculated. When $D$ represents one intended
population, its prevalence can remain implicit. Reporting AP at the
evaluation's own prevalence answers a question about that sample. Resampling to
a target ratio discards data while still fixing a single population. Reporting
the precision--recall curve defers the choice to a reader who lacks the
sampling information needed to make it. Each leaves the claim indexed to one
population. A first task may therefore be to construct a target-indexed
quantity, $\AP_\pi(D)$, that allows for calculation of average precision at any
chosen prevalence using all of the available data. Such a quantity carries the
selection rates observed in the evaluation over to the target population
unchanged, which is a claim about that population rather than a property of the
arithmetic.

One target-prevalence value does not render a performance claim applicable to
several different deployment populations. Any single convenient prevalence
hides where performance weakens, while averaging over a distribution of target
prevalences both requires that distribution to be defensible and permits a
strong result at one prevalence to offset a weak result elsewhere. The claim
pursued here is the more demanding one: that a performance level holds
throughout a scientifically plausible prevalence range rather than in
expectation across it.

A second problem arises if the claim is intended to persist under new
evaluation units or under other sources of variation. This calls for an
evaluation regime that specifies which empirical tests fall within the scope
of the claim. The extent to which a claim survives those attempts at
refutation is what will be called its support.

The population and replication problems cannot be solved independently.
Different evaluations may weaken at different prevalences, so challenging every
evaluation at one common prevalence can miss the condition that defeats each of
them. A claim required to hold throughout the population set in every
evaluation must allow the limiting prevalence to change between evaluations.

Neither problem concerns the precision of an estimate. Variation under a
declared computational resampling law describes sensitivity to recombinations
of the observed units; it does not say whether the advantage would remain in a
population with a different supply of negatives, or in the next evaluation.
What is at issue is survival rather than precision: whether the advantage
remains once the conditions under which it could fail have been named and
imposed.

Call $D$ a full evaluation when it contains the outcomes and scores needed to
evaluate every declared condition on one set of units, so that no condition
requires a separate dataset. Let
\[
\Assess=(\CSet,\Regime,K),
\]
where $\CSet$ contains the conditions that can be evaluated within each full
evaluation, $\Regime$ specifies the admissible full evaluations, and $K$ is
the order of the replication challenge. For AP, target prevalence will supply
the condition set. Changes whose effects cannot be evaluated by varying
prevalence must be represented by distinct evaluations admitted by $\Regime$.
The order $K$ is the number of distinct evaluation effects in each subset
used by the replication challenge. It is neither estimated from data nor
equipped with a default. It states how severe a test the claimant has chosen
to face, while the number of observed evaluations states how much empirical
evidence is available. An observed order-$K$ assessment therefore requires at
least $K$ separately conducted evaluations. If fewer are available, that
assessment is unavailable; $K$ should not be lowered after the results are
seen merely to obtain a verdict.
These elements belong to the claim rather than to the analysis. Like a
preregistered analysis plan, they lose their force if they are chosen after the
comparison has been seen.

Probability measures introduced below are counterfactual or computational
characterizations. They are not asserted to be data-generating laws for the
observed units or for unseen evaluations. The support construction instead
acts directly on a finite indexed effect list. A target distribution states a
declared population description, and a resampling law defines a computational
projection.

What can then be assessed depends on the evidence available. Resampling a
single evaluation supplies a computational projection under a declared
resampling procedure. It can recombine only the case types, scores, and failure
mechanisms represented in that evaluation. The resulting support is projected rather than observed, and no
number of resamples converts one evaluation into new empirical tests.

The motivating application we develop is model replacement: whether a candidate
model $A$ should replace an incumbent model $B$ on ranking performance. Pairing
introduces a third difficulty. Subtracting two separately challenged reports
lets each model meet its adverse condition in a different population or a
different evaluation. Forming the difference on the same units removes that
latitude at the cost of a new one, since finite-sample AP can reward whichever
model divides tied scores more finely, whether or not the division tracks the
outcomes.


\section{Average precision across target prevalences}
\label{sec:prevalence}

\subsection{Prior-standardized average precision}

Changing the observed positive fraction can alter both the population represented and the evidence used to describe the ranking. Prior standardization separates these operations by asking a counterfactual question: what precision would the observed class-conditional ranking behavior produce under a different class prior?

Within one evaluation, let $S_j$ be the score produced by the model instance used for $j\in\{A,B\}$. That instance and its scoring rule remain fixed throughout prior standardization. Writing $S$ generically for $S_A$ or $S_B$, let the observed class-conditional selection rates at threshold $\tau$ be
\[
r(\tau)=\Prb_{\mathrm{obs}}(S\geq \tau\mid Y=1),
\qquad
f(\tau)=\Prb_{\mathrm{obs}}(S\geq \tau\mid Y=0).
\]
Here $\Prb_{\mathrm{obs}}$ characterizes the observed score--outcome
distribution, while $\Prb_{\mathrm{target}}$ below characterizes a declared
target population. Neither notation posits a stochastic law that generated the
evaluation data.
Holding these rates fixed while varying the target prevalence is warranted only when the class-conditional score behavior transports. Because the two models are compared on the same units, the transport condition for each target population must preserve their joint class-conditional score behavior:
\begin{equation}
\Prb_{\mathrm{obs}}\{(S_A,S_B)\in G\mid Y=y\}
=
\Prb_{\mathrm{target}}\{(S_A,S_B)\in G\mid Y=y\},
\qquad y\in\{0,1\},
\label{eq:score-invariance}
\end{equation}
for every measurable $G\subseteq\mathbb R^2$. This condition entails the corresponding marginal invariance for $S_A$ and $S_B$, which supplies each model's target AP profile. It also preserves the within-class association between the scores in paired computational projections. This score-level form of prior-probability shift is weaker than invariance of $P(X\mid Y)$ but follows from it for fixed scoring rules \cite{saerens2002adjusting,storkey2009training,lipton2018detecting}.

Under this condition, Bayes' rule expresses precision at threshold $\tau$ in a target population with prevalence $\pi\in(0,1)$ as
\begin{equation}
p_\pi(\tau)
=
\frac{\pi r(\tau)}
{\pi r(\tau)+(1-\pi)f(\tau)}.
\label{eq:reference-precision}
\end{equation}
Its odds separate the target prior from the ranking behavior:
\begin{equation}
\frac{p_\pi(\tau)}{1-p_\pi(\tau)}
=
\frac{\pi}{1-\pi}\frac{r(\tau)}{f(\tau)}.
\label{eq:precision-odds}
\end{equation}
The functions $r(\tau)$ and $f(\tau)$ are therefore held at the values supplied by evaluation $D$ while $\pi$ indexes the declared target populations. Varying $\pi$ changes only the class prior. Aggregating $p_\pi(\tau)$ across the distinct score thresholds using the observed positive-class recall increments gives $\AP_\pi(D)$, prior-standardized AP for evaluation $D$ at prevalence $\pi$ \cite{elkan2001foundations,siblini2020calibration}. Appendix~\ref{app:ap} gives the empirical definition.

The class-conditional score-transport condition matters because prevalence can accompany other population changes. A referral population may contain more severe cases within each outcome class, while a new measurement process may change the score for otherwise similar cases. In either setting, Equation~\eqref{eq:reference-precision} reweights rates that need not describe the target. Those changes must be represented by observed evaluations from the corresponding environments rather than by prior adjustment alone \cite{bareinboim2013transportability}.

\subsection{A common effect scale}

Prior-standardized AP still has a different random-ranking value at every target prevalence. Population random ranking has $r(\tau)=f(\tau)$ and hence $\AP_\pi=\pi$. Define
\begin{equation}
\CNAP_\pi(D)
=
\frac{\AP_\pi(D)-\pi}{1-\pi}.
\label{eq:cnap}
\end{equation}
CNAP is zero for population random ranking, one for perfect ranking, and negative below chance. The normalization has the chance-corrected form $(\text{observed}-\text{chance})/(\text{ceiling}-\text{chance})$ \cite{cohen1960coefficient}. For two models,
\begin{equation}
\CNAP_{A,\pi}-\CNAP_{B,\pi}
=
\frac{\AP_{A,\pi}-\AP_{B,\pi}}{1-\pi}.
\label{eq:cnap-difference}
\end{equation}
The common chance term cancels, while division by the random-to-perfect headroom $1-\pi$ expresses the AP gap on a common opportunity scale. Thus a margin $d$ requires a raw AP gain of $d(1-\pi)$ at prevalence $\pi$; the same identity applies after tie averaging. This scaling is substantive and is appropriate when the replacement margin refers to a fraction of random-to-perfect improvement rather than an equal raw AP increment at every prevalence.

Common endpoints do not make the profile prevalence-invariant. At a threshold with true-positive and false-positive rates $r_h$ and $f_h$, the normalized contribution is
\begin{equation}
\frac{p_{\pi,h}-\pi}{1-\pi}
=
\frac{\pi(r_h-f_h)}
{\pi r_h+(1-\pi)f_h}.
\label{eq:threshold-cnap}
\end{equation}
Rare-positive populations therefore charge false positives more heavily. Two models can exchange order as $\pi$ changes even when both profiles improve with prevalence.

Let $\PSet\subset(0,1)$ be a nonempty compact set chosen from deployment plans or scientific knowledge before the paired-difference profile is examined. The CNAP profile then supplies the AP-specific condition index: $c=\pi$ and $\CSet=\PSet$. No prevalence has yet been selected or averaged away. The next problem is to determine what a claim throughout this profile must retain when the evaluation is replicated.

\section{Support across declared conditions and evaluations}
\label{sec:support}

The prevalence profile in Section~\ref{sec:prevalence} is one instance of a condition-indexed effect. More generally, let $\Delta_c(D)$ denote an effect under condition $c\in\CSet$, oriented so that larger values favor the claim. Every condition in the nonempty set $\CSet$ must be evaluable from one full evaluation $D$. For AP under score-level prior invariance, $c$ is a target prevalence and $\CSet=\PSet$. Section~\ref{sec:paired} constructs the paired, tie-averaged CNAP difference that supplies $\Delta_\pi(D)$ for model replacement.

Two sources of challenge enter through different empirical objects. Conditions in $\CSet$ vary within each full evaluation. Changes that alter class-conditional ranking behavior require full evaluations conducted in the corresponding contexts. The evaluation regime states which of those evaluations address the declared claim.

\subsection{The evaluation regime}

Let $\Regime$ specify the requirements under which a full evaluation $D$ counts as an admissible empirical test of the claim. These requirements include the object under test, eligible populations, measurement standards, available outcomes and scores, and the independent evaluation unit. The regime defines scientific scope. It makes no assertion that admissible evaluations share a statistical distribution. Statistical descriptions applied after observation may differ between contexts.

A separately conducted full evaluation satisfying $\Regime$ is an \emph{empirical replication}. Replications use nonoverlapping evaluation units. Their model construction and data collection are not altered in response to another evaluation in the same replication set. Observations within an evaluation may be clustered or matched when the design and analysis preserve that structure. Administrative partitions, resampling draws, and cross-validation folds do not create empirical replications.

A fixed-model regime holds the fitted models and measurement process fixed while admitting new evaluation units. A broader regime may admit distinct environments or a prespecified model-updating procedure. When models are updated, the common object under test is the procedure that produces them; information used for updating must remain separate from the subsequent evaluation. In that setting, $A$ and $B$ denote procedures whose context-specific instances supply the scores in each evaluation. If an earlier evaluation prompts a later update, the evaluations form an adaptive evidential sequence. They do not constitute a set of separate empirical replications under the definition above, and their order and information flow require explicit analysis.

Several empirical replications may occur in different domain contexts while remaining admissible tests of the same scoped claim. Their identity as separate tests should be retained when contextual variation is scientifically meaningful. If dataset membership is only administrative and all units address one aggregate population, pooling may be more appropriate. Scientific reach is determined by $\Regime$ and by the empirical tests conducted within it.

\subsection{The effect retained across all conditions}

Let $D_1,\ldots,D_L$ be a nonempty finite list of full evaluations admitted by $\Regime$. A claim throughout $\CSet$ requires one effect level to hold under every declared condition within each evaluation. For evaluation $D_\ell$, the greatest such level is
\begin{equation}
z_\ell
=
\inf_{c\in\CSet}\Delta_c(D_\ell).
\label{eq:evaluation-retained-effect}
\end{equation}
This is the largest uniform lower bound on the condition profile: every declared condition retains at least $z_\ell$, and no larger level has that property. The infimum covers finite condition sets and continuous profiles whose lowest value may fail to be attained. An average across conditions would require weights over $\CSet$ and would permit favorable conditions to compensate for unfavorable ones.

This operation is algebraically related to an intersection bound. Its role here
is different: Equation~\eqref{eq:evaluation-retained-effect} is evaluated
exactly on the condition profile supplied by one finite evaluation, rather than
used to estimate an infimum of population-indexed functions. Section~\ref{sec:related-frameworks}
develops this distinction.

Suppose three evaluations have the profiles
\[
\begin{array}{c@{\qquad}rr@{\qquad}r}
\toprule
& \pi=0.01 & \pi=0.40 & z_\ell \\
\midrule
D_1 & -0.10 & \phantom{-}0.30 & -0.10 \\
D_2 & \phantom{-}0.20 & -0.20 & -0.20 \\
D_3 & \phantom{-}0.30 & \phantom{-}0.30 & \phantom{-}0.30 \\
\bottomrule
\end{array}
\]
At the two common prevalences, the mean effects are $0.133\overline{3}$ and
$0.133\overline{3}$. Equation~\eqref{eq:evaluation-retained-effect} gives
$z_1=-0.10$, $z_2=-0.20$, and $z_3=0.30$ because the limiting prevalence
differs between the first two evaluations.

Collect the retained effects in $\mathbf z=(z_1,\ldots,z_L)$. The replication challenge determines how strongly low values in this finite list influence the assessment.

\subsection{Support across retained evaluation effects}

The empirical mean of $\mathbf z$ permits a large retained effect in one evaluation to offset a small or negative retained effect in another. Its minimum gives the level retained by every available evaluation. The order $K$ supplies graded lower-tail emphasis between these endpoints.

\begin{definition}[Replication-survival support]
\label{def:support}
For $1\leq K\leq L$, let
\[
\mathcal I_{L,K}
=
\{I\subseteq\{1,\ldots,L\}:|I|=K\}.
\]
Define
\begin{equation}
S_{K,\CSet}(\Delta;\mathbf D)
=
S_K(\mathbf z)
=
\binom{L}{K}^{-1}
\sum_{I\in\mathcal I_{L,K}}
\min_{\ell\in I}z_\ell.
\label{eq:robust-support}
\end{equation}
\end{definition}

$S_K(\mathbf z)$ is the average weakest effect over all subsets of $K$
distinct evaluation indices. Each subset receives equal weight, and no
evaluation is repeated within a challenge. Equal numerical effects can still
represent distinct empirical tests. This construction supplies no probability
law for unseen evaluations. A different weighting would define a different
assessment and would require a scientific basis.

At $K=1$, support is the empirical mean of the retained effects. At $K=2$,
\begin{equation}
S_2(\mathbf z)
=
\overline z
-
\frac{1}{L(L-1)}
\sum_{1\leq i<\ell\leq L}|z_i-z_\ell|.
\label{eq:k2}
\end{equation}
The second-order criterion is the mean retained effect minus half the mean
pairwise disagreement over distinct ordered pairs. As $K$ increases, lower
retained effects receive more weight. The endpoints are
$S_1(\mathbf z)=\overline z$ and $S_L(\mathbf z)=\min_\ell z_\ell$.

\begin{samepage}
\noindent For the retained-effect list $(0.70,0.35,0.05)$,
$S_1=0.366\overline{6}$. The order-2 construction contains three equally
weighted subsets:
\[
\{0.70,0.35\},\qquad
\{0.70,0.05\},\qquad
\{0.35,0.05\}.
\]
Their minima are $0.35$, $0.05$, and $0.05$, giving
\[
S_2
=
\frac{0.35+0.05+0.05}{3}
=0.15.
\]
The sole order-3 subset contains all three effects, so $S_3=0.05$.
\end{samepage}
Thus the sequence $0.366\overline{6},0.15,0.05$ moves from the mean to
the observed minimum. The list length records the three available effects;
$K$ controls how many distinct effects must jointly face each challenge.

For any nonempty finite effect list $\mathbf t=(t_1,\ldots,t_L)$ and
$1\leq K\leq L$, $S_K(\mathbf t)$ denotes the same distinct-subset
construction with $\mathbf t$ in place of $\mathbf z$.

\subsection{Literal survival at a declared floor}

Supported magnitude retains the numerical size of every subset minimum. A
decision may also require the retained effect to clear a prespecified floor
$d$. Let
\[
R_d
=
\sum_{\ell=1}^L\mathbf 1\{z_\ell>d\}.
\]
Empirical literal survival is
\begin{equation}
\psi_{K,d}(\mathbf z)
=
\frac{\binom{R_d}{K}}{\binom{L}{K}},
\label{eq:literal-survival-main}
\end{equation}
where $\binom{R_d}{K}=0$ when $R_d<K$. It is the fraction of distinct
order-$K$ subsets whose effects all clear $d$.

For any level $s$, define
\[
R_s=\sum_{\ell=1}^L\mathbf 1\{z_\ell>s\},
\qquad
z_-=\min_\ell z_\ell,
\qquad
z_+=\max_\ell z_\ell.
\]
Then
\begin{equation}
\begin{aligned}
S_K(\mathbf z)
&=
z_-
+
\int_{z_-}^{z_+}
\psi_{K,s}(\mathbf z)\,ds \\
&=
z_-
+
\int_{z_-}^{z_+}
\frac{\binom{R_s}{K}}{\binom{L}{K}}\,ds.
\end{aligned}
\label{eq:support-survival-main}
\end{equation}
At each level $s$, the integrand is the fraction of distinct order-$K$
subsets whose effects all clear $s$. Supported magnitude integrates this
finite-list survival fraction over all effect levels. Literal survival
evaluates it at $d$. For a generic finite list $\mathbf t$, the notation
$\psi_{K,d}(\mathbf t)$ applies the same construction.
Appendix~\ref{app:literal-survival} gives the ordered representation and
further details.

\subsection{Observed and projected effect lists}
\label{sec:projected}

Let $N$ count the evaluation units whose outcomes and scores enter an AP calculation, $M$ count separately conducted empirical evaluations, and $J$ count computational replications.

The generic list length $L$ in the preceding definitions can then be
instantiated as either $M$ or $J$ depending on the nature of the evaluations.
For example, for projected effects from $J$ computational replications of one
observed evaluation, $L=J$ and the list index identifies the replications. In
either case, the order must satisfy $K\leq L$.

When $M$ separately conducted full evaluations satisfy $\Regime$, their retained effects form
\[
\mathbf z_M^{\mathrm{obs}}
=
(z_1,\ldots,z_M).
\]
When $K\leq M$, applying Equations~\eqref{eq:robust-support} and
\eqref{eq:literal-survival-main} directly to this indexed list gives observed
support and literal survival. If $M<K$, the declared observed assessment is
unavailable because the evidence does not contain enough distinct evaluations
to form an order-$K$ subset.

One observed evaluation has not empirically faced several replications. Resampling it generates a list of $J$ computational effects under a declared resampling procedure:
\[
\mathbf z_J^{\mathrm{proj}}
=
(z_1,\ldots,z_J).
\]
When $J\geq K$, applying the same finite-list formulas gives
\emph{projected support} and projected literal survival. These quantities
characterize the effects produced by the procedure from the observed
evaluation. The procedure uses only score values and failure modes represented
there, so its challenge is narrower than one supplied by new empirical
replications. In practice, $J$ is usually chosen much larger than $K$ to
control Monte Carlo error.

Separate empirical replications can reveal failures arising from variation within the scope of $\Regime$ that the original evaluation did not contain. An evaluation involving a mechanism excluded from $\Regime$ may show that the declared assessment was too narrow.

Increasing $J$ reduces Monte Carlo error in the same projected quantities. Increasing $K$ changes the challenge applied to either effect list. Increasing $M$ changes the empirical evidence available for observed support.

\subsection{Challenge order and scope}

Equation~\eqref{eq:robust-support} first identifies the limiting condition within each evaluation and then applies the replication challenge to the retained effects. The condition attaining $z_\ell$ may therefore differ between evaluations.

For a fixed condition $c$, write $\Delta_c[\mathbf D]=(\Delta_c(D_1),\ldots,\Delta_c(D_L))$ for the effect list obtained by applying that condition to every evaluation. Reversing the order gives
\[
\inf_{c\in\CSet}
S_K\left(\Delta_c[\mathbf D]\right)
\]
and requires one condition to challenge all evaluations before the limiting
condition is selected. In the profiles displayed above, $D_1$ is weakest at
$\pi=0.01$, $D_2$ is weakest at $\pi=0.40$, and $D_3$ has the same effect at
both prevalences. At $K=2$, the retained-effect order gives
\[
S_2(-0.10,-0.20,0.30)=-\frac16.
\]
The common-condition values are
$S_2(-0.10,0.20,0.30)=0$ at $\pi=0.01$ and
$S_2(0.30,-0.20,0.30)=-1/30$ at $\pi=0.40$, giving $-1/30$ after the
infimum. The retained-effect order supplies the stronger test within the
declared assessment: each evaluation faces every declared condition, and the
condition exposing its weakest effect may differ between evaluations.

The retained-effect order always satisfies
\begin{equation}
S_K(\mathbf z)
\leq
\inf_{c\in\CSet}
S_K\left(\Delta_c[\mathbf D]\right).
\label{eq:order}
\end{equation}
Equality holds when the same condition minimizes every relevant profile; Appendix~\ref{app:support-proofs} gives the proof.

Enlarging $\CSet$ gives each evaluation more conditions under which its effect may be limited. Increasing $K$ gives greater influence to the lower retained effects.

\begin{proposition}[Challenge monotonicity]
\label{prop:monotonicity}
If $\mathcal C_1\subseteq\mathcal C_2$ and
$1\leq K_1\leq K_2\leq L$, then
\[
S_{K_2,\mathcal C_2}(\Delta;\mathbf D)
\leq
S_{K_1,\mathcal C_1}(\Delta;\mathbf D).
\]
\end{proposition}

$\Regime$ determines which evaluations and procedures are admissible. Arbitrary changes of $\Regime$ have no comparable ordering because they change the scientific scope of the assessment.

Support has a limited negative interpretation. A level is retained when it
survives a challenge constructed from the declared conditions, the admitted
empirical tests, and the chosen distinct-subset order. Survival establishes no
truth beyond that assessment \cite{popper1959logic,popper1963conjectures}. The
regime supplies admissibility, the conducted evaluations supply the finite
indexed effect list, and $K$ sets the challenge applied to that list.

\section{Constructing the paired AP effect}
\label{sec:paired}

The support operator requires a finite list of evaluation-level effects. For replacement, each effect compares candidate $A$ with incumbent $B$ within the same evaluation and target prevalence. Its AP calculation should give no expected credit to outcome-blind refinement of score ties.

\subsection{Marginal summaries do not compose}

A replacement decision concerns an advantage realized under the same conditions. For $\mathbf D=(D_1,\ldots,D_L)$, define
\[
\begin{aligned}
a_\ell
&=
\inf_{\pi\in\PSet}\CNAP_{A,\pi}(D_\ell),
&\qquad
b_\ell
&=
\inf_{\pi\in\PSet}\CNAP_{B,\pi}(D_\ell),
\\
t_\ell
&=
\inf_{\pi\in\PSet}
\{\CNAP_{A,\pi}(D_\ell)-\CNAP_{B,\pi}(D_\ell)\}.
\end{aligned}
\]
Let $\mathbf a=(a_1,\ldots,a_L)$, with $\mathbf b$ and $\mathbf t$ defined similarly. In $S_K(\mathbf a)-S_K(\mathbf b)$, the adverse prevalence may differ between models within an evaluation, and the two support calculations may emphasize different evaluations. The discrepancy has a fixed direction:
\begin{equation}
S_K(\mathbf a)
-
S_K(\mathbf b)
\geq
S_K(\mathbf t).
\label{eq:noncomposition}
\end{equation}
Separate reports therefore overstate the supported paired advantage by an amount they do not identify. The proof in Appendix~\ref{app:paired-proofs} uses only properties of the infimum and minimum, so the same result applies after tie averaging is introduced below.

\subsection{Ordinary pairing rewards arbitrary refinement}

Pairing aligns the units and outcomes, but ordinary finite-sample AP still depends on the score vectors' tie structures. Stepwise AP is upward displaced under random ordering because it evaluates precision when positive cases enter the ranking \cite{bestgen2015random}. A fully tied score vector has one threshold and returns the population random-ranking value, while a vector of distinct scores creates more opportunities for favorable positive entries.

Four units show why this matters for replacement. Let model $A$ assign four distinct scores and model $B$ assign one common score. With two positive labels, two negative labels, and $\pi=1/2$, assign the labels uniformly over the six possibilities. Neither score vector is associated with the outcomes. Under the ordinary threshold-block convention, model $A$ has expected AP $49/72$, while model $B$ has AP $1/2$ under every assignment. Hence
\[
\E[\CNAP_A-\CNAP_B]
=
\frac{13}{36}.
\]
For a $K=2$ challenge under two independently randomized label assignments, the expected minimum of the paired differences remains positive at $29/216$. Pairing has favored the model that broke ties without information.

No single additive reference correction removes this displacement at every level of model performance. The expected paired statistic under joint label permutation is large near no association and vanishes at perfect separation, where both models reach the ceiling. Appendix~\ref{app:paired-proofs} gives an exact four-unit example in which subtracting the joint-permutation expectation favors the coarser model even though both models separate the classes perfectly. Tie handling must therefore be resolved before the nonlinear population and replication challenges are applied.

\subsection{Tie-averaged average precision}
\label{sec:tie-average}

Let $\mathcal O(D)$ be the strict rankings consistent with a score vector, obtained by ordering tied units in every possible way while preserving the order of the score blocks.

\begin{definition}[Tie-averaged average precision]
\begin{equation}
\APta_\pi(D)
=
\frac{1}{|\mathcal O(D)|}
\sum_{o\in\mathcal O(D)}
\AP_\pi(D_o),
\label{eq:tie-average}
\end{equation}
where $D_o$ is evaluation $D$ resolved into strict ranking $o$. Define $\CNAPta_\pi$ through Equation~\eqref{eq:cnap}.
\end{definition}

The convention corresponds to a declared random or otherwise outcome-blind resolution of units within each score block. It is appropriate when an operational cutoff may split a block and the model supplies no within-block information. If deployment selects whole score blocks or uses a prespecified informative secondary rule, threshold-block AP or the resulting ordered ranking may be more faithful. Under the random-resolution regime, tie averaging is equivalent to adding independent continuous jitter within blocks and averaging ordinary AP; Appendix~\ref{app:tie-computation} evaluates that average without enumeration or random jitter.

Tie averaging still rewards informative refinement. A finer score improves realized AP when it orders cases correctly within a coarse block and loses AP when it orders them incorrectly. Only refinement unrelated to the outcomes receives zero expected credit.

\begin{proposition}[Neutrality to arbitrary refinement]
\label{prop:refinement}
Suppose model $A$ preserves the order of model $B$'s score blocks and, conditionally on $S_B$ and $Y$, its strict ranking is uniform over the rankings consistent with $S_B$. Then, for every $\pi$,
\[
\E\!\left[
\APta_{A,\pi}\mid S_B,Y
\right]
=
\APta_{B,\pi}.
\]
\end{proposition}

The proposition conditions on the realized outcomes, so it applies when the coarse model is informative. It removes expected credit for arbitrary refinement without requiring exchangeable evaluation units.

\subsection{The supported paired advantage}

The AP-specific condition index is the target prevalence. Define the paired effect profile
\[
\Delta^{A:B}_\pi(D)
=
\CNAPta_{A,\pi}(D)
-
\CNAPta_{B,\pi}(D),
\]
and let $\mathbf D=(D_1,\ldots,D_L)$. The retained effects from Section~\ref{sec:support} are
\[
z_\ell^{A:B}
=
\inf_{\pi\in\PSet}\Delta^{A:B}_\pi(D_\ell),
\qquad
\mathbf z^{A:B}
=
(z_1^{A:B},\ldots,z_L^{A:B}).
\]
The model-replacement support is the corresponding instance of Equation~\eqref{eq:robust-support}:
\begin{equation}
\theta_{A:B;K}(\mathbf D)
=
S_{K,\PSet}\!\left(\Delta^{A:B};\mathbf D\right)
=
S_K\!\left(\mathbf z^{A:B}\right).
\label{eq:paired-estimand}
\end{equation}
Within each evaluation, the construction subtracts tie-averaged CNAP at each prevalence and then takes the prevalence infimum. It collects the retained advantages and applies the replication challenge. A positive value of $\theta_{A:B;K}(\mathbf D)$ means that the average weakest paired advantage over subsets of $K$ distinct evaluation effects is positive. It describes retained magnitude; literal survival evaluates the same effect list at a declared floor.

\subsection{Orientation and no verdict}

The reverse effect profile is $\Delta^{B:A}_\pi(D)=-\Delta^{A:B}_\pi(D)$. Define
\[
z_\ell^{B:A}
=
\inf_{\pi\in\PSet}\Delta^{B:A}_\pi(D_\ell),
\qquad
\mathbf z^{B:A}
=
(z_1^{B:A},\ldots,z_L^{B:A}),
\]
and
\[
\theta_{B:A;K}(\mathbf D)
=
S_K\!\left(\mathbf z^{B:A}\right).
\]
Negation occurs before the prevalence infimum and the replication challenge, so the reversed value is generally different from the negative of the original. The two orientations satisfy
\begin{equation}
\theta_{A:B;K}(\mathbf D)+\theta_{B:A;K}(\mathbf D)\leq0.
\label{eq:orientation}
\end{equation}
The construction cannot support positive advantages in both directions.

Both orientations may be nonpositive because the declared challenges defeat superiority in either direction. The diagnostic range from $\theta_{A:B;K}(\mathbf D)$ to $-\theta_{B:A;K}(\mathbf D)$ reflects variation across target prevalences and disagreement across evaluations; it is not an uncertainty interval. Appendix~\ref{app:paired-proofs} gives its characterization.

\subsection{Reference behavior and design scope}
\label{sec:paired-scope}

Refinement neutrality compares a fine ranking with a coarser ranking it refines. Reference behavior for two arbitrary uninformative rankings is evaluated by a constructed diagnostic randomization.

\begin{definition}[Exchangeable-label reference law]
\label{def:exchangeable}
Conditionally on both score vectors and the positive count $N_+$, every label vector containing $N_+$ positives has probability $\binom{N}{N_+}^{-1}$.
\end{definition}

Under this law, labels read along any fixed strict ranking are uniform over the same set of arrangements. Expected tie-averaged AP therefore depends on $N_+$, $N_-$, and $\pi$, but not on the score vector or its tie structure.

Write $\E_0$ for expectation when the exchangeable-label randomization is applied independently within each full evaluation, conditional on the score vectors and positive counts fixed by the design. This expectation is a diagnostic reference value, or anchor, showing whether score structure alone shifts the comparison under the randomization.

\begin{proposition}[Reference behavior of paired support]
\label{prop:paired-reference}
Under the exchangeable-label reference law,
\[
\E_0[\Delta^{A:B}_\pi(D)]=0
\qquad\text{for every }\pi\in\PSet,
\]
and consequently
\[
\E_0[\theta_{A:B;K}(\mathbf D)]\leq0,
\qquad
\E_0[\theta_{B:A;K}(\mathbf D)]\leq0.
\]
\end{proposition}

The inequality need not be equality because the prevalence infimum and replication minima are nonlinear. They can give an uninformative paired effect a negative expected supported value. The result prevents differing tie structures from manufacturing a positive supported advantage, although the challenge itself may be conservative.

Full exchangeability is stronger than independence. Matched, clustered, and heterogeneous-risk designs usually permit only a restricted group of label arrangements, under which score vectors aligned with the strata can have different reference values. Tie averaging still removes arbitrary ordering within score blocks, and Equation~\eqref{eq:paired-estimand} remains descriptive, but Proposition~\ref{prop:paired-reference} no longer follows. Appendix~\ref{app:paired-proofs} gives a matched-pair counterexample.

\subsection{A two-part replacement rule}

For a prespecified floor $d\geq0$, define literal survival of the paired advantage as
\begin{equation}
\psi^{A:B}_{K,d}(\mathbf D)
=
\psi_{K,d}\!\left(\mathbf z^{A:B}\right).
\label{eq:paired-literal-survival}
\end{equation}
Both quantities summarize the same replication-challenged paired effect.
Supported magnitude retains the numerical size of every subset minimum.
Literal survival records the fraction of distinct order-$K$ subsets whose
effects all exceed $d$.

Let $\delta\geq0$ be the smallest replication-challenged CNAP advantage that would justify replacement, and let $\gamma\in(0,1)$ be the minimum acceptable literal-survival fraction. The two-part policy requires
\begin{equation}
\theta_{A:B;K}(\mathbf D)>\delta
\qquad\text{and}\qquad
\psi^{A:B}_{K,d}(\mathbf D)>\gamma.
\label{eq:replacement-targets}
\end{equation}
The magnitude condition prevents a negligible but broadly retained advantage from automatically justifying replacement and remains sensitive to severe evaluation effects. The survival condition prevents large favorable effects from compensating for an unacceptable empirical fraction of failure at $d$. Separating $\delta$ from $d$ clarifies these questions: for example, $\delta>0$ and $d=0$ require a worthwhile retained advantage together with a high empirical survival fraction for positive superiority. Setting $\delta=d$ is permissible when the same practical effect level is intended for both purposes.

The severity of the literal-survival requirement depends jointly on $K$ and $\gamma$. If $R_d$ of $M$ empirical effects clear $d$, observed literal survival is
\begin{equation}
\psi_{K,M,d}^{A:B,\mathrm{obs}}
=
\frac{\binom{R_d}{K}}{\binom{M}{K}}.
\label{eq:empirical-survival-fraction}
\end{equation}
For $J$ computational effects, projected literal survival replaces $M$ by
$J$. The fraction is zero whenever fewer than $K$ effects clear $d$ and is one
only when every effect clears $d$. Reports should give $R_d/M$ for observed
support or $R_d/J$ for projected support, together with the corresponding
subset survival, $K$, and $\gamma$.

Define $\psi^{B:A}_{K,d}(\mathbf D)=\psi_{K,d}(\mathbf z^{B:A})$. The reverse decision applies the same two-part rule to $\theta_{B:A;K}(\mathbf D)$ and $\psi^{B:A}_{K,d}(\mathbf D)$, with orientation-specific thresholds if the consequences are asymmetric. If neither orientation satisfies both requirements, the analysis returns no verdict.

Equation~\eqref{eq:replacement-targets} is a finite-evidence policy rule. It
applies to the observed evaluation effects or to a declared computational
projection, as identified in the report. Passing the rule supports replacement
relative to those evaluations and challenges; it is not a confidence statement
about a latent population quantity or about unseen evaluations. Failure of
either requirement returns no verdict rather than evidence that the incumbent
is superior.

The rule does not infer that future deployment will reproduce the observed
effects. Using its verdict to guide deployment requires an external decision
premise: that the declared condition set, evaluation regime, and transport
commitments are adequate for the contemplated use, together with acceptance of
the consequences of residual failure. Those commitments belong to the claimant
or decision maker. The support calculation makes them explicit but neither
supplies nor certifies them.

\section{Evidence and computational projection}
\label{sec:evidence}

Evidence is supplied by observed full evaluations. Resampling an evaluation
transforms that evidence into computational effects under a declared procedure;
it does not add evidence. The same support functionals apply to observed and
computational effect lists, while their interpretations differ. When $M$
separately conducted evaluations satisfying $\Regime$ are available, their
effects give observed support through Equation~\eqref{eq:robust-support};
Appendix~\ref{app:empirical-evaluations} gives the calculation.

\subsection{Projected support from one observed evaluation}

Most applications begin with one observed full evaluation. Its relationship to model development and to any earlier evaluations must be reported. A declared empirical-resampling procedure can project the support calculation within the empirical world represented by that evaluation.

Let $D_{\mathrm{obs},+}$ and $D_{\mathrm{obs},-}$ contain the positive and negative evaluation units. Computational replication $j$ samples the declared numbers of units with replacement within each class, keeps both model scores attached to every sampled unit, and evaluates the paired effect profile throughout $\PSet$. Its retained effect is
\[
\begin{aligned}
z_j^{A:B}
&=
\inf_{\pi\in\PSet}\Delta^{A:B}_\pi(D^{(j)}) \\
&=
\inf_{\pi\in\PSet}
\{\CNAPta_{A,\pi}(D^{(j)})-\CNAPta_{B,\pi}(D^{(j)})\}.
\end{aligned}
\]
The computational effect list
\[
\mathbf z_J^{A:B,\mathrm{proj}}
=
(z_1^{A:B},\ldots,z_J^{A:B})
\]
is an indexed list of finite computational effects. Its projected supported
magnitude is
\[
\widehat\theta_{A:B;K,J}^{\,\mathrm{proj}}
=
S_K\!\left(\mathbf z_J^{A:B,\mathrm{proj}}\right).
\]

The same computational replications give
\begin{equation}
\widehat\psi_{K,J,d}^{A:B,\mathrm{proj}}
=
\frac{\binom{R_d^{(J)}}{K}}{\binom{J}{K}},
\qquad
R_d^{(J)}=\sum_{j=1}^J\mathbf 1\{z_j^{A:B}>d\}.
\label{eq:projected-literal-survival}
\end{equation}
The empirical resampling law is a constructed computational measure on
recombinations of the observed units. The exact projected quantities are the
corresponding support functionals under that law. The displayed quantities use
$J$ draws to approximate them. Increasing $J$ refines the same computational
projection and reduces its Monte Carlo variation. It supplies no additional
empirical opportunity to reveal failure in a new context.

Stratification fixes the positive and negative counts in the projected calculation. A projection of the observed evaluation uses its class counts. A prospective projection uses the counts planned for a later evaluation. A procedure that varies the class counts must declare how they are resampled. Clustered or matched studies must resample the declared design blocks rather than their component observations, which may preclude simple class-stratified resampling.

The same resampled units must be used for both models and throughout the prevalence profile. Resampling models separately would destroy pairing. Resampling separately at each prevalence would splice together values from different computational evaluations, so their minimum would no longer describe one evaluation's entire prevalence profile.

The empirical-resampling projection represents only recombinations available in the observed score pools. It cannot produce unseen score values or failures caused by a new environment, measurement process, or development run. Its result must be reported as projected fixed-model support. It records no survival of $K$ empirical replications.

\subsection{Lower-tail sensitivity and sparse outcomes}

The role of $K$ creates a computational and interpretive tension. After sorting $J$ projected computational effects, Equation~\eqref{eq:robust-support} assigns $z_{(i)}$ weight
\[
w_{i,K,J}
=
\frac{\binom{J-i}{K-1}}{\binom{J}{K}},
\qquad i=1,\ldots,J-K+1,
\]
with zero weight on later order statistics. For $M$ observed empirical
effects, the same formula replaces $J$ by $M$.
Increasing $K$ moves weight toward the lower observed effects. The stronger challenge consequently depends more heavily on whether the empirical evaluations, or the resampling collection, contain adverse evaluation effects. This dependence is especially consequential when positive outcomes are sparse and AP is coarse. Literal survival introduces an additional boundary sensitivity: when several $z_\ell$ lie near $d$, small perturbations can change the exceedance indicator even when the effect changes little.

\subsection{Prevalence search and nonsmooth minima}

A finite set $\PSet=\{\pi_1,\ldots,\pi_G\}$ defines an exact finite challenge. A grid used to approximate an interval can miss an interior minimum and raise the reported value. A one-dimensional optimizer can instead search a compact interval, with its tolerance checked against a dense grid.

Nonunique or moving prevalence minimizers require a further distinction. Uniformly accurate profiles may still yield an accurate infimum value even when the minimizing prevalence is unstable. Because the infimum is nonsmooth at flat or multiple minima, the reported minimum and the identity of the limiting prevalence can respond differently to small numerical or data perturbations. Reports should therefore separate accuracy of the minimum value from stability of its location.

\subsection{Computational uncertainty}

Observed support is an exact functional of the finite observed effect list, up
to numerical evaluation of the condition infimum. Projected support adds
Monte Carlo approximation because only $J$ draws from the declared empirical
resampling law are used. Neither source is sampling uncertainty about a latent
population parameter.

Computational diagnostics should address the numerical tolerance of the
prevalence search, stability across random seeds, convergence as $J$ increases,
and sensitivity when effects lie near $d$. Pairing, tie averaging, and the same
prevalence search must be preserved in every computational replication. Monte
Carlo error may be summarized because the randomization is deliberately
introduced by the algorithm; it quantifies approximation of the declared
projection and supplies no claim about unseen units, contexts, or evaluations.

\section{Interpretation, reporting, and scope}
\label{sec:discussion}

\subsection{What the quantities support}

Supported magnitude and literal survival depend on the declared assessment and
on the evaluations supplied to it; neither is a portable property of either
model. They describe the same retained paired effects. Supported magnitude
averages the weakest effect over subsets of $K$ distinct evaluation effects,
while literal survival records the fraction of those subsets whose effects all
clear $d$. In the AP application, $\CSet=\PSet$ and the effect is the
tie-averaged CNAP difference. Omitting any of those declarations changes what
the quantities mean.

Projected and observed summaries require different language. Several separately conducted empirical replications give observed support across the tests satisfying the declared regime. A one-evaluation empirical-resampling procedure gives projected magnitude and survival by recombining the observed within-class score pools and should not be described as having survived external replications.

The result also remains a ranking-performance claim. A deployment decision may turn on calibration or on the consequences of acting on the selected cases, neither of which Equation~\eqref{eq:replacement-targets} evaluates.

\subsection{Relation to population-based robustness and inference}
\label{sec:related-frameworks}

Several constructions in this paper have analogues in literatures that begin
with a population probability law. The shared algebra does not imply a shared
interpretation. Here support acts directly on subsets of a finite indexed list
of retained effects. It neither posits nor estimates a superpopulation law for
unseen evaluations. Probability enters only when a declared target population
or deliberately constructed computational experiment requires it.

The retained effect $z_\ell=\inf_{c\in\CSet}\Delta_c(D_\ell)$ resembles an
intersection bound. Chernozhukov, Lee, and Rosen study infima and suprema of
population-indexed functions and develop bias correction and repeated-sampling
inference for their bounds \cite{chernozhukov2013intersection}. Equation~\eqref{eq:evaluation-retained-effect}
instead extracts the exact uniform lower level of the finite condition profile
computed in evaluation $D_\ell$. It neither estimates a functional
$\inf_c\theta_c(P)$ nor receives a confidence interpretation from
intersection-bound asymptotics.

Replication-survival support also has a finite spectral form. Let
$\ell_1,\ldots,\ell_L$ be the losses $-z_1,\ldots,-z_L$, with ordered values
$\ell_{(1)}\leq\cdots\leq\ell_{(L)}$. Then
\begin{equation}
-S_K(\mathbf z)
=
\sum_{r=K}^L
\frac{\binom{r-1}{K-1}}{\binom{L}{K}}\,
\ell_{(r)}.
\label{eq:spectral-loss-relation}
\end{equation}
The nondecreasing weights give $-S_K$ a finite spectral analogy
\cite{acerbi2002spectral}, but they depend on both $L$ and $K$. The manuscript
invokes neither the associated coherence axioms nor a population-risk
interpretation. Equation~\eqref{eq:support-survival-main} is likewise related
algebraically to distortion and Choquet representations of survival functions
\cite{wang1996premium}. In those literatures the distribution is ordinarily
the law of a random quantity and empirical order statistics estimate a
population functional. Here Equations~\eqref{eq:spectral-loss-relation} and
\eqref{eq:quantile-weight} are exact representations of a finite indexed list.

The capped AUROC weight sets in Appendix~\ref{app:auroc} are similarly related
to distributionally robust optimization, which optimizes over alternative
probability measures near a reference measure \cite{shapiro2017distributionally}.
Ambiguity sets in that literature may be selected through divergence geometry
and calibrated as confidence regions for uncertain probabilities
\cite{bental2013robust}. The present cap instead declares how strongly the
represented positive and negative cases may be redistributed. The factor
$\Gamma$ is a deterministic concentration allowance, not a confidence radius,
a model of population shift, or a claim that an unknown true distribution lies
in the admissible set.

Consequently, population consistency, asymptotic normality, and coverage
results from these adjacent literatures are not invoked. Data-driven robust
optimization can obtain such coverage by assuming independent observations
from an unknown population law \cite{duchi2021statistics}; that is a different
inferential target. Randomness deliberately introduced by the computational
procedure in Section~\ref{sec:evidence} supports Monte Carlo statements only
about approximation of the declared projection. Additional empirical reach
requires additional evaluations within $\Regime$, not a stronger probabilistic
interpretation of the evaluations already observed.

\subsection{Minimum reporting content}

A concise report should identify the claim before presenting diagnostics. Table~\ref{tab:reporting} gives the minimum content.

\begin{table}[ht]
\centering
\caption{Minimum reporting content for a supported replacement analysis.}
\label{tab:reporting}
\begin{tabularx}{\textwidth}{@{}p{0.23\textwidth}>{\raggedright\arraybackslash}X@{}}
\toprule
Component & Report \\
\midrule
Claim and decision &
Model orientation; observed $\theta_{A:B;K,M}^{\mathrm{obs}}$ and $\psi_{K,M,d}^{A:B,\mathrm{obs}}$, or projected $\widehat\theta_{A:B;K,J}^{\mathrm{proj}}$ and $\widehat\psi_{K,J,d}^{A:B,\mathrm{proj}}$; the reversed orientation; $\delta$, $d$, and $\gamma$; and the resulting finite-evidence verdict \\
Assessment scope &
The admissibility requirements in $\Regime$; $\CSet$ and its scientific basis, with $\CSet=\PSet$ for AP; what is fixed or may differ; the independent evaluation unit; and $K$ \\
Evidence &
The separately conducted empirical evaluations supplying the evidence; whether a computational projection was subsequently constructed; $N$ or the evaluation-specific $N_m$; $M$ or $J$; confirmation that the list length is at least $K$; observed class counts; evaluation contexts; the number clearing $d$ and the resulting fraction of distinct order-$K$ subsets; and the empirical-resampling procedure where applicable \\
Effect construction &
The invariance used to evaluate every $c\in\CSet$ from one evaluation; paired evaluation and resampling; the tie-averaged convention; and whether the diagnostic exchangeable-label randomization applies, with its scientific basis \\
Computation and diagnostics &
Condition set or optimizer; numerical tolerance; computational-replication count $J$, Monte Carlo error, seed stability, and convergence where applicable; paired-effect profiles and minimizing conditions; evaluation effects; values near $d$; and sensitivity to $K$, $\CSet$, and prespecified alternative floors \\
\bottomrule
\end{tabularx}
\end{table}

A possible headline statement for one observed evaluation is:
\begin{quote}
Under the declared assessment, model $A$ was compared with model $B$ using paired tie-averaged CNAP throughout $\PSet$. The analysis used $J$ computational replications from one observed evaluation and reports projected support. Supported magnitude was [value], and literal survival above $d=[value]$ was [value]. The reversed orientation was [values]. Monte Carlo error was [value], and the prevalence search used tolerance [value]. With $\delta=[value]$ and $\gamma=[value]$, the finite-evidence rule returned [supported replacement/no verdict] relative to the observed evaluation and declared computational projection.
\end{quote}
With several empirical evaluations, the second sentence should state, ``The analysis used $M$ separately conducted full evaluations satisfying the declared regime and reports observed support.'' The report should give the individual contexts and the number of evaluations that cleared $d$.

\subsection{Limits of the AP claim}

Prior standardization does not transport ranking behavior across changes within outcome class. If severity, measurement, or feature distributions change $P(S\mid Y)$, target-population evaluations are required.

Tie averaging does not create full label exchangeability. In matched, clustered, or heterogeneous-risk designs, the paired quantity remains computable but may inherit a design-specific anchor unless the admissible comparisons respect the design's symmetry.

The choices of $\PSet$, $\Regime$, and $K$ remain scientific commitments. Widening $\PSet$ or increasing $K$ mechanically lowers the criterion, but neither operation proves that the resulting assessment captures the failures relevant to deployment.

\subsection{Beyond prevalence: empirical case-mix breakdown for AUROC}
\label{sec:auroc-outlook}

The support construction also applies beyond prevalence and AP. Pure prior shift leaves AUROC unchanged, while changes in the mix of cases within each outcome class can challenge an AUROC advantage. Appendix~\ref{app:auroc} defines an empirical concentration challenge and the case-mix breakdown factor at which the replacement policy first fails. The result remains limited to the case types represented in the evaluation.

\section{Conclusion}

A model-replacement claim based on average precision must survive more than a favorable marginal summary. The population must be named because AP changes with prevalence, while the evaluations must be named because an average advantage can conceal a failure to persist. These are instances of the same noncompensation requirement, but they enter the assessment differently: target prevalence can be evaluated within each replication, whereas other variation must be represented by separate evaluations admitted by $\Regime$.

Prior standardization evaluates each model at the target prevalences when class-conditional score behavior is invariant. Chance normalization places those values on a common population random-to-perfect scale. Tie averaging then removes expected credit for arbitrary refinement before the models are compared on the same units. For $J$ computational replications from one observed evaluation, the prevalence infimum is applied within each replication before projected order-$K$ support:
\[
\widehat\theta_{A:B;K,J}^{\mathrm{proj}}
=
\binom{J}{K}^{-1}
\sum_{\substack{I\subseteq\{1,\ldots,J\}\\|I|=K}}
\min_{j\in I}
\inf_{\pi\in\PSet}
\{\CNAPta_{A,\pi}(D^{(j)})-\CNAPta_{B,\pi}(D^{(j)})\}.
\]
For $M$ separately conducted empirical evaluations, observed support has the same form with $M$ and $D_m$.

Supported magnitude and empirical literal survival summarize different features
of the same challenged advantage. The finite-evidence replacement rule requires
retained magnitude to exceed $\delta$ and the fraction of distinct order-$K$
subsets whose effects all clear floor $d$ to exceed $\gamma$. The first
requirement preserves practical magnitude and sensitivity to severe losses;
the second prevents favorable outcomes from compensating for excessive failure
at the declared floor. Passing the rule supports replacement relative to the
observed evaluations and declared challenges, not as a confidence statement
about a latent population quantity.

Several empirical replications give both quantities from the same evaluation effects. One evaluation can support only empirical-resampling projections under a declared computational procedure. Those projections can guide a replacement analysis, but they cannot substitute computational resampling for new empirical opportunities to reveal failure.

The condition-indexed support operator also applies beyond prevalence. For AUROC, a symmetric capped concentration challenge can search over increasingly adverse redistributions of the observed positive and negative cases and report the empirical case-mix breakdown factor at which the same replacement policy first fails. This extension supplies a reproducible empirical robustness assessment when no target mix is selected. The search rule also removes the need to choose one concentration bound for the primary analysis. The assessment remains silent about case types absent from the evaluation.

\appendix

\section{Empirical prior-standardized average precision}
\label{app:ap}

\subsection{Definition and weighted implementation}

Let
\[
D=\{(s_i,y_i)\}_{i=1}^{N},
\qquad
y_i\in\{0,1\},
\]
where $N$ is the number of evaluation units on which AP is calculated. Let $N_+=\sum_i y_i>0$ and $N_-=N-N_+>0$ be the positive and negative counts. Let $c_1>\cdots>c_H$ be the distinct observed score values. At threshold $h$, define
\[
\operatorname{TP}_h
=
\sum_i\mathbf 1\{y_i=1,s_i\geq c_h\},
\qquad
\operatorname{FP}_h
=
\sum_i\mathbf 1\{y_i=0,s_i\geq c_h\},
\]
and
\[
r_h=\frac{\operatorname{TP}_h}{N_+},
\qquad
f_h=\frac{\operatorname{FP}_h}{N_-}.
\]
With $r_0=0$, noninterpolated prior-standardized AP under the threshold-block convention is
\begin{equation}
\AP_\pi(D)
=
\sum_{h=1}^{H}
(r_h-r_{h-1})
\frac{\pi r_h}
{\pi r_h+(1-\pi)f_h}.
\label{eq:empirical-ap}
\end{equation}
All units sharing a score enter one threshold block. A block containing only negative cases has zero recall increment, although its false positives reduce the precision of later blocks.

The same quantity can be computed by weighting each positive unit by $w_+=\pi/N_+$ and each negative unit by $w_-=(1-\pi)/N_-$. Weighted precision at threshold $h$ is Equation~\eqref{eq:reference-precision}, while weighted recall remains $r_h$.

\subsection{Behavior across prevalence}

Combining Equations~\eqref{eq:cnap} and \eqref{eq:empirical-ap} gives
\begin{equation}
\CNAP_\pi(D)
=
\sum_{h=1}^{H}
(r_h-r_{h-1})
\frac{\pi(r_h-f_h)}
{\pi r_h+(1-\pi)f_h}.
\label{eq:cnap-profile}
\end{equation}
The derivative of a threshold contribution is
\begin{equation}
\frac{\partial}{\partial\pi}
\left\{
\frac{\pi(r_h-f_h)}
{\pi r_h+(1-\pi)f_h}
\right\}
=
\frac{f_h(r_h-f_h)}
{\{\pi r_h+(1-\pi)f_h\}^2}.
\label{eq:cnap-derivative}
\end{equation}
If $r_h\geq f_h$ at every threshold with positive recall increment, the profile is nondecreasing and its minimum over $[\pi_L,\pi_U]$ occurs at $\pi_L$. A paired difference can remain nonmonotone even when both marginal profiles are nondecreasing.

\section{Properties of replication-survival support}
\label{app:support-proofs}

\subsection{Survival and ordered representations}

For a finite list $\mathbf t=(t_1,\ldots,t_L)$ with entries in $[a,b]$, let
\[
R_s=\sum_{\ell=1}^L\mathbf 1\{t_\ell>s\}.
\]
For each $I\in\mathcal I_{L,K}$,
\[
\min_{\ell\in I}t_\ell
=
a+\int_a^b\mathbf 1\{t_\ell>s\text{ for every }\ell\in I\}\,ds.
\]
Among the $\binom{L}{K}$ subsets, exactly $\binom{R_s}{K}$ have every
effect above $s$. Averaging the preceding identity over the subsets gives
\begin{equation}
S_K(\mathbf t)
=
a+\int_a^b
\frac{\binom{R_s}{K}}{\binom{L}{K}}\,ds.
\label{eq:survival-representation}
\end{equation}
At each level $s$, the integrand is the fraction of distinct order-$K$
subsets for which every effect retains the level.

\subsection{Literal survival of a practical margin}
\label{app:literal-survival}

For a prespecified practical margin $d$, define
\begin{equation}
\psi_{K,d}(\mathbf t)
=
\frac{\binom{R_d}{K}}{\binom{L}{K}}.
\label{eq:literal-survival}
\end{equation}
This quantity is a finite-list subset fraction at one effect level.
Equation~\eqref{eq:survival-representation} can be written
\[
S_K(\mathbf t)=a+\int_a^b\psi_{K,s}(\mathbf t)\,ds,
\]
so replication-survival support integrates finite-list literal survival over
all effect levels.

For an observed effect list, literal survival is an exact fraction of its
distinct subsets. For a projected list, finite $J$ introduces only Monte Carlo
approximation to the declared empirical-resampling projection. Because AP is
discrete, the choice between $>$ and $\geq$ and sensitivity to values near $d$
must be declared.

Let $t_{(1)}\leq\cdots\leq t_{(L)}$. The minimum of a distinct order-$K$
subset equals $t_{(i)}$ for exactly $\binom{L-i}{K-1}$ subsets. Hence
\begin{equation}
S_K(\mathbf t)
=
\binom{L}{K}^{-1}
\sum_{i=1}^{L-K+1}
\binom{L-i}{K-1}t_{(i)}.
\label{eq:quantile-weight}
\end{equation}
The weights move toward the lower observed effects as $K$ increases. This exact
finite-list representation has the spectral and distortion analogues discussed
in Section~\ref{sec:related-frameworks} \cite{acerbi2002spectral,wang1996premium}.

\subsection{Proof of the operator-order inequality}

For each $I\in\mathcal I_{L,K}$, let
$X_c^{(\ell)}=\Delta_c(D_\ell)$ and
$Y_{c,I}=\min_{\ell\in I}X_c^{(\ell)}$. Then
\[
\min_{\ell\in I}\inf_{c\in\CSet}X_c^{(\ell)}
=
\inf_{c\in\CSet}\min_{\ell\in I}X_c^{(\ell)}
=
\inf_{c\in\CSet}Y_{c,I}.
\]
Since $\inf_c Y_{c,I}\leq Y_{c',I}$ for every $c'\in\CSet$, averaging
over $I\in\mathcal I_{L,K}$ gives
\[
\binom{L}{K}^{-1}\sum_{I\in\mathcal I_{L,K}}
\inf_{c\in\CSet}Y_{c,I}
\leq
\inf_{c\in\CSet}
\binom{L}{K}^{-1}\sum_{I\in\mathcal I_{L,K}}Y_{c,I},
\]
which proves Equation~\eqref{eq:order}.

\subsection{Proof of challenge monotonicity}

Take $\mathcal C_1\subseteq\mathcal C_2$ and
$1\leq K_1\leq K_2\leq L$. Expanding the condition set lowers each
evaluation effect, so it remains to compare the subset orders for one fixed
list. Every order-$K_2$ subset contains $\binom{K_2}{K_1}$ order-$K_1$
subsets, and its minimum is no greater than the average minimum of those
subsets. Averaging this inequality over all order-$K_2$ subsets counts each
order-$K_1$ subset exactly $\binom{L-K_1}{K_2-K_1}$ times. The resulting
averages are $S_{K_2}$ and $S_{K_1}$, which proves
Proposition~\ref{prop:monotonicity}.

\section{Paired comparison proofs and counterexamples}
\label{app:paired-proofs}

\subsection{Why marginal subtraction overstates paired advantage}

Let $U_\pi$ and $V_\pi$ be two profiles, with
\[
U_{\PSet}=\inf_{\pi\in\PSet}U_\pi,
\qquad
V_{\PSet}=\inf_{\pi\in\PSet}V_\pi,
\qquad
T_{\PSet}=\inf_{\pi\in\PSet}(U_\pi-V_\pi).
\]
Choosing a prevalence minimizing $U_\pi$ gives
\[
T_{\PSet}
\leq
U_{\PSet}-V_\pi
\leq
U_{\PSet}-V_{\PSet}.
\]
Monotonicity of $S_K$ therefore gives
\[
S_K(U_{\PSet}-V_{\PSet})
\geq
S_K(T_{\PSet}).
\]

For effects evaluated on the same replications, the minimum is superadditive:
\[
\min_j(X_j+Y_j)
\geq
\min_jX_j+\min_jY_j.
\]
Taking $X=U_{\PSet}-V_{\PSet}$ and $Y=V_{\PSet}$ yields
\[
S_K(U_{\PSet})-S_K(V_{\PSet})
\geq
S_K(U_{\PSet}-V_{\PSet}),
\]
which completes Equation~\eqref{eq:noncomposition}.

\subsection{Proof of refinement neutrality}

Condition on $S_B$ and $Y$. Proposition~\ref{prop:refinement} makes model $A$'s strict ranking $O_A$ uniform on $\mathcal O(D_B)$. Tie-averaged AP for $A$ is the conditional expectation of ordinary AP over whatever ties $A$ retains. By the tower property,
\[
\E[\APta_{A,\pi}\mid S_B,Y]
=
\E[\AP_\pi(D_{O_A})\mid S_B,Y].
\]
Uniformity of $O_A$ turns the right side into the average of $\AP_\pi$ over $\mathcal O(D_B)$, which is $\APta_{B,\pi}$.

\subsection{Rank invariance under exchangeable labels}

Fix $N_+$, $N_-$, and a strict ranking $o$. Under Definition~\ref{def:exchangeable}, the label sequence read along $o$ is uniform over the $\binom N{N_+}$ arrangements containing $N_+$ positives. Reading along another strict ranking is a bijection of the same arrangements, so the sequence has the same law.

For a strict ranking, $\pi\mapsto\CNAP_\pi(D_o)$ is a deterministic function of that sequence. Its expectation is therefore the same for every strict ranking. Tie-averaged AP is a convex combination of strict-ranking values, so any two score vectors on the same units satisfy
\[
\E_0[\CNAPta_{A,\pi}]
=
\E_0[\CNAPta_{B,\pi}].
\]
Thus $\E_0[\Delta^{A:B}_\pi(D)]=0$.

For a fixed evaluation $D$,
\[
\E_0\!\left[
\inf_{\pi\in\PSet}\Delta^{A:B}_\pi(D)
\right]
\leq
\inf_{\pi\in\PSet}\E_0[\Delta^{A:B}_\pi(D)]
=0,
\]
and $S_K(\mathbf t)\leq L^{-1}\sum_{\ell=1}^Lt_\ell$ for any finite list of length $L$. Applying this inequality to every label randomization and then taking $\E_0$ proves Proposition~\ref{prop:paired-reference}; exchanging $A$ and $B$ proves the reversed orientation.

\subsection{Why restricted label symmetry is insufficient}

Take four units in two matched pairs, each containing exactly one positive. The design permits labels to swap within pairs but not between pairs. Let model $A$ give all units the same score, while model $B$ places the first pair above the second and ties the units within each pair. At $\pi=1/2$, exact tie averaging over the four permitted assignments gives
\[
\CNAPta_A-\CNAPta_B
=
\frac1{36}
\]
in one orientation for every assignment. The effect has no replication
variation, so its supported value remains $1/36$ for every admissible $K$.

The obstruction is the symmetry group of the design. Full exchangeability places every strict ranking in one orbit of the label-permutation group. Restricted designs may divide rankings into several orbits with different anchors, and averaging within score ties cannot move a ranking between them.

\subsection{Why paired null centering fails}

Take two positive and two negative units at $\pi=1/2$, with
\[
s_A=(4,3,2,1),
\qquad
s_B=(2,2,1,1),
\]
and let the two highest-scored units be positive. Both models separate the classes perfectly, so every stratified empirical-resampling evaluation has CNAP one for each model and paired difference zero.

Under joint label permutation, exact enumeration of the $K=2$ support calculation gives an expected supported value of $35/384$ for $A-B$ and $-295/1152$ for $B-A$. Subtracting these reference values would report a positive corrected value of about $0.256$ for the second orientation, favoring the coarser model although both predictions are perfect. The displacement is not an additive constant because perfect separation leaves no room for it.

\subsection{Proof of the orientation bound}

For evaluation $j$, write
\[
L_j=\inf_{\pi\in\PSet}\Delta^{A:B}_\pi(D_j),
\qquad
U_j=\sup_{\pi\in\PSet}\Delta^{A:B}_\pi(D_j).
\]
For each $I\in\mathcal I_{L,K}$, the two orientations contribute
\[
\min_{j\in I}L_j
\qquad\text{and}\qquad
-\max_{j\in I}U_j.
\]
Since $L_j\leq U_j$ for every evaluation,
\[
\min_{j\in I}L_j-\max_{j\in I}U_j
\leq0,
\]
and averaging over the distinct order-$K$ subsets proves
Equation~\eqref{eq:orientation}. The negative of this sum is the width of the
no-verdict region. When $\PSet$ is a singleton and $K=2$, the width is
\[
\binom{L}{2}^{-1}\sum_{1\leq i<j\leq L}|T_i-T_j|
\]
for the finite paired-effect list.

\section{Closed-form computation of tie-averaged AP}
\label{app:tie-computation}

Index the $H$ distinct score blocks by $h=1,\ldots,H$ from highest to lowest. Let block $h$ contain $a_h$ positive and $b_h$ negative units, with $m_h=a_h+b_h$, and define
\[
A_h=\sum_{\ell<h}a_\ell,
\qquad
B_h=\sum_{\ell<h}b_\ell.
\]
Then
\begin{equation}
\APta_\pi(D)
=
\sum_{h=1}^{H}
\frac{a_h}{N_+}\frac1{m_h}
\sum_{k=1}^{m_h}
\sum_{t=\max(0,k-1-b_h)}^{\min(k-1,a_h-1)}
\frac{\binom{a_h-1}{t}\binom{b_h}{k-1-t}}
{\binom{m_h-1}{k-1}}
\frac{\pi r_{hkt}}
{\pi r_{hkt}+(1-\pi)f_{hkt}},
\label{eq:tie-closed}
\end{equation}
where
\[
r_{hkt}=\frac{A_h+t+1}{N_+},
\qquad
f_{hkt}=\frac{B_h+k-1-t}{N_-}.
\]

To obtain the formula, fix one positive unit in block $h$. Under a uniformly random ordering of the block, its position $k$ is uniform on $\{1,\ldots,m_h\}$. Given $k$, the number $t$ of other positive units in the $k-1$ preceding positions is hypergeometric with $a_h-1$ available positives and $b_h$ negatives. The final fraction in Equation~\eqref{eq:tie-closed} is the prior-standardized precision when that positive enters the ranking. Averaging over $t$ and $k$, multiplying by the block's $a_h$ positive units, and summing over blocks gives the result.

The calculation costs $O(\sum_hm_h^2)$ operations. Its combinatorial weights do not depend on $\pi$ and can be reused throughout the prevalence search. When all scores are distinct, Equation~\eqref{eq:tie-closed} reduces to ordinary stepwise AP.

\section{Computation and approximation}
\label{app:computation}

\subsection{Several empirical evaluations}
\label{app:empirical-evaluations}

Suppose $M$ separately conducted full evaluations $D_1,\ldots,D_M$ satisfy $\Regime$. Let $N_m$ be the number of evaluation units in $D_m$. The values $N_1,\ldots,N_M$ may differ, and the evaluations may occur in distinct admissible contexts. For each evaluation, compute
\[
z_m
=
\inf_{c\in\CSet}\Delta_c(D_m).
\]
Let $\mathbf z_M^{\mathrm{obs}}=(z_1,\ldots,z_M)$. When $K\leq M$,
observed support is
\begin{equation}
S_{K,M,\CSet}^{\,\mathrm{obs}}
=
S_K(\mathbf z_M^{\mathrm{obs}})
=
\binom{M}{K}^{-1}
\sum_{I\in\mathcal I_{M,K}}
\min_{m\in I}z_m.
\label{eq:observed-empirical-support}
\end{equation}
For AP replacement, $\CSet=\PSet$ and the effect profile is $\Delta^{A:B}_\pi$, giving $\theta_{A:B;K,M}^{\,\mathrm{obs}}$. The same effects give
\begin{equation}
\psi_{K,M,d}^{A:B,\mathrm{obs}}
=
\frac{\binom{R_d}{K}}{\binom{M}{K}},
\qquad
R_d=\sum_{m=1}^M\mathbf 1\{z_m>d\}.
\label{eq:literal-survival-estimator}
\end{equation}
These are exact functionals of the finite indexed effect list. Each distinct
evaluation receives equal influence in the subset construction. A small $M$
gives few distinct subsets. Increasing $K$ shifts support toward the lower
effects and requires more effects to clear $d$ jointly. If $M<K$, the declared
observed assessment is unavailable. The count $M$ records the number of
observed evaluations.

If the fitted models change between evaluations, the evaluations address one scoped claim only when a common prespecified procedure produces those models. The count $M$ includes only evaluations whose model construction was not altered in response to another evaluation in the set. Updates prompted by an earlier evaluation create an adaptive sequence that requires a separate ordered analysis. Dataset partitions with no scientific identity should not be used to increase $M$.

\subsection{Sorted computation}

Sort the $M$ observed evaluation effects:
\[
z_{(1)}\leq\cdots\leq z_{(M)}.
\]
Exactly $\binom{M-i}{K-1}$ distinct order-$K$ subsets have minimum
$z_{(i)}$. Hence Equation~\eqref{eq:observed-empirical-support} equals
\begin{equation}
S_{K,M,\CSet}^{\,\mathrm{obs}}
=
\binom{M}{K}^{-1}
\sum_{i=1}^{M-K+1}
\binom{M-i}{K-1}z_{(i)}.
\label{eq:sorted-observed}
\end{equation}

For $J$ empirical-resampling effects, the Monte Carlo approximation is
\begin{equation}
\widehat S_{K,J,\CSet}^{\,\mathrm{proj}}
=
\binom{J}{K}^{-1}
\sum_{i=1}^{J-K+1}
\binom{J-i}{K-1}z_{(i)}.
\label{eq:sorted-bootstrap}
\end{equation}
At $K=2$, this reduces to
\[
\widehat S_{2,J,\CSet}^{\,\mathrm{proj}}
=
\frac{2}{J(J-1)}
\sum_{i=1}^{J-1}(J-i)z_{(i)}.
\]

\subsection{Monte Carlo error}

For $K=2$, define the leave-one-replication projection
\[
\widehat h_{1,-i}(z_i)
=
\frac1{J-1}\sum_{\substack{j=1\\j\ne i}}^{J}\min(z_i,z_j).
\]
A first-order Monte Carlo standard error for the projected approximation is
\begin{equation}
\widehat{\operatorname{se}}_{\mathrm{MC}}
=
\frac2{\sqrt J}
\operatorname{sd}_{i}
\left\{
\widehat h_{1,-i}(z_i)
\right\}.
\label{eq:mcse}
\end{equation}
Prefix sums over the sorted array compute all projections in linear time after
sorting.

\subsection{Prevalence evaluation and numerical checks}

For a finite prevalence set with $G$ values, score-block counts are computed once in each evaluation and reused across prevalences. For a continuous interval, Equation~\eqref{eq:cnap-profile} is a smooth one-dimensional objective away from thresholds with no recall contribution. Endpoint values and stationary points may be checked directly, while a numerical optimizer should be validated against a dense grid. The report should give the optimizer tolerance and should distinguish stability of the minimum value from stability of the minimizing prevalence. For projected support, calculations across larger values of $J$ and independent random seeds assess Monte Carlo convergence to the declared empirical-resampling projection.

\section{Empirical case-mix breakdown for AUROC}
\label{app:auroc}

\subsection{Prior shift supplies no AUROC challenge}

The condition set used for AP is a set of target prevalences. It is available
because AP depends on the class prior. AUROC does not. AUROC is the probability
that a randomly drawn positive outranks a randomly drawn negative, with half
credit for a tie, and is therefore a functional of the two class-conditional
score distributions alone. Changing the prior alters how often each class is
drawn while leaving both conditionals intact. An AUROC advantage that survives
every $\pi\in\PSet$ has consequently survived no additional challenge.

The within-class composition remains free to change, however. A new population
may contain different proportions of mild and severe positives or easy and
difficult negatives even when its prevalence is unchanged. Named target mixes
can be evaluated when those proportions and the required transport assumptions
are available. For settings without named target mixes, the observed cases
define the case support for the challenge. The procedure admits increasingly
concentrated redistributions within each outcome class and reports the first
concentration at which either replacement criterion fails. It begins with
uniform within-class weights and progressively relaxes a common upper bound on
the weight any observation may receive. Reporting the breakdown factor removes
the need to select one concentration bound for the primary analysis. The bound
and the resulting breakdown factor are defined below.

\subsection{The paired AUROC difference as a case-mix average}

Let $D$ contain paired scores $(s_i^A,s_i^B)$ and outcome $Y_i$ for every
evaluation unit, with positive index set $\mathcal I_+$, negative index set
$\mathcal I_-$, and class counts $N_+$ and $N_-$. Thus the dataset contains
$N=N_++N_-$ evaluation units. For model $m\in\{A,B\}$,
write
\[
\phi_m(i,k)
=
\mathbf 1\{s_i^m>s_k^m\}
+\tfrac12\mathbf 1\{s_i^m=s_k^m\},
\qquad i\in\mathcal I_+,\quad k\in\mathcal I_-.
\]
Thus $\phi_m(i,k)$ is one when model $m$ orders the positive--negative
pair correctly, one half for a tie, and zero otherwise. Define the pairwise
gain matrix
\[
G_{ik}=\phi_A(i,k)-\phi_B(i,k).
\]
If positive and negative cases receive normalized weights $w_+$ and $w_-$,
the paired AUROC difference under that empirical case mix is
\begin{equation}
\Delta_D(w_+,w_-)
=
\sum_{i\in\mathcal I_+}\sum_{k\in\mathcal I_-}
w_{+,i}w_{-,k}G_{ik}.
\label{eq:auroc-weighted-difference}
\end{equation}
Uniform within-class weights recover the ordinary empirical AUROC difference.
Because both models are evaluated with the same $w_+$ and $w_-$, the
redistribution remains paired.

No additional tie correction is needed. Each pair already receives half credit
for a tie, so exchangeable labels give every fixed score vector mean AUROC one
half regardless of its tie structure. The positive finite-sample reference
value that required tie-averaged AP in Section~\ref{sec:tie-average} does not
recur.

\subsection{Why unrestricted redistribution degenerates}

If the weights in Equation~\eqref{eq:auroc-weighted-difference} are otherwise
unrestricted, all positive weight may be placed on one positive case and all
negative weight on one negative case. The infimum then reduces to the least
favorable single entry of $G$. This is too concentrated to support a useful
replacement criterion.

The degeneracy is especially clear when candidate $A$ separates perfectly.
Then $G_{ik}=1-\phi_B(i,k)\geq0$ for every pair, but concentrating on any pair
that incumbent $B$ also orders correctly gives difference zero. A perfect
candidate can show a strictly positive worst-case advantage only if the
incumbent fails on every pair. Unrestricted empirical reweighting therefore
does not make replacement appropriately demanding; it generally makes a
strict advantage unavailable.

\subsection{The symmetric capped concentration challenge}

What causes the degeneracy is concentration rather than empirical reweighting
itself. For a common factor $\Gamma\geq1$, define
\begin{equation}
\mathcal W_{+,\Gamma}(D)
=
\left\{
w_+:
\sum_{i\in\mathcal I_+}w_{+,i}=1,\quad
0\leq w_{+,i}\leq\frac{\Gamma}{N_+}
\right\},
\label{eq:auroc-positive-weights}
\end{equation}
and
\begin{equation}
\mathcal W_{-,\Gamma}(D)
=
\left\{
w_-:
\sum_{k\in\mathcal I_-}w_{-,k}=1,\quad
0\leq w_{-,k}\leq\frac{\Gamma}{N_-}
\right\}.
\label{eq:auroc-negative-weights}
\end{equation}
The declared condition set is
\[
\CSet_\Gamma(D)
=
\mathcal W_{+,\Gamma}(D)\times\mathcal W_{-,\Gamma}(D).
\]
This capped density family is a finite empirical analogue of an ambiguity set
in distributionally robust optimization \cite{shapiro2017distributionally,bental2013robust}.
Here it is fixed as a deterministic challenge over represented cases; it is not
calibrated as a confidence region for an unknown distribution.
The advantage retained in evaluation $D$ is
\begin{equation}
Z_\Gamma^{A:B}(D)
=
\inf_{(w_+,w_-)\in\CSet_\Gamma(D)}
\Delta_D(w_+,w_-).
\label{eq:auroc-z-gamma}
\end{equation}

At $\Gamma=1$, normalization forces every weight to be uniform, so
\begin{equation}
Z_1^{A:B}(D)
=
\operatorname{AUC}_A(D)-\operatorname{AUC}_B(D).
\label{eq:auroc-gamma-one}
\end{equation}
At larger values, no observation may receive more than $\Gamma$ times its
uniform within-class weight. The bound is one-sided: an observation may receive
zero weight, while concentration on the retained observations is limited from
above. The effective sample size (ESS) summarizes this concentration. For
either class $y$,
\[
\sum_i w_{y,i}^2
\leq
\left(\max_i w_{y,i}\right)\sum_iw_{y,i}
\leq
\frac{\Gamma}{N_y},
\]
and hence
\begin{equation}
\operatorname{ESS}(w_y)
=
\left(\sum_iw_{y,i}^2\right)^{-1}
\geq
\frac{N_y}{\Gamma}.
\label{eq:auroc-gamma-ess}
\end{equation}
Thus $1/\Gamma$ is a guaranteed minimum effective-support fraction implied by
the cap, although an effective-sample-size restriction alone would not define
the same set of weights.

The admissible sets are nested:
\[
\Gamma_1\leq\Gamma_2
\quad\Longrightarrow\quad
\CSet_{\Gamma_1}(D)\subseteq\CSet_{\Gamma_2}(D),
\]
so
\begin{equation}
Z_{\Gamma_2}^{A:B}(D)
\leq
Z_{\Gamma_1}^{A:B}(D).
\label{eq:auroc-gamma-monotone}
\end{equation}
A perfectly separating model continues to attain AUROC one throughout every
admissible set. The challenge controls the maximum concentration allowed among
the observed cases in each outcome class.

The common factor $\Gamma$ gives both classes the same normalized upper
allowance. Their realized redistributions remain separate because the positive
and negative weights are optimized independently. Suppose separate bounds
$\Gamma_+$ and $\Gamma_-$ were introduced. Every combination satisfying
$1\leq\Gamma_+\leq\Gamma$ and $1\leq\Gamma_-\leq\Gamma$ would lie in a square
of weaker concentration allowances. Monotonicity makes $(\Gamma,\Gamma)$ its
most adverse corner. Equation~\eqref{eq:auroc-z-gamma} therefore challenges
that square with one concentration coordinate. Asymmetric bounds may be useful
as diagnostics. The primary implementation and minimum report use the common
factor.

\subsection{Replication-challenged decision profiles}

Apply the paper's support construction separately at every common concentration
factor. For a finite list $\mathbf D$ of empirical or computational evaluations, let
\[
\mathbf z(\Gamma;\mathbf D)
=
Z_\Gamma^{A:B}[\mathbf D].
\]
Define the supported-magnitude profile
\begin{equation}
\theta_{A:B;K}(\Gamma;\mathbf D)
=
S_K\!\left(\mathbf z(\Gamma;\mathbf D)\right)
\label{eq:auroc-theta-gamma}
\end{equation}
and, for an AUROC-difference floor $d\geq0$, the literal-survival profile
\begin{equation}
\psi_{K,d}^{A:B}(\Gamma;\mathbf D)
=
\psi_{K,d}\!\left(\mathbf z(\Gamma;\mathbf D)\right).
\label{eq:auroc-psi-gamma}
\end{equation}
The same common $\Gamma$ is applied to every evaluation, while the adverse
weights may differ between evaluations. Because
Equation~\eqref{eq:auroc-gamma-monotone} holds pointwise, both
$\theta_{A:B;K}(\Gamma;\mathbf D)$ and $\psi_{K,d}^{A:B}(\Gamma;\mathbf D)$ are nonincreasing in
$\Gamma$.

Let $\delta\geq0$ now denote the smallest supported AUROC advantage that would
justify replacement, and retain $\gamma\in(0,1)$ as the required literal
survival fraction. At concentration factor $\Gamma$, the AUROC replacement
policy requires
\begin{equation}
\theta_{A:B;K}(\Gamma;\mathbf D)>\delta
\qquad\text{and}\qquad
\psi_{K,d}^{A:B}(\Gamma;\mathbf D)>\gamma.
\label{eq:auroc-policy-gamma}
\end{equation}
At $\Gamma=1$, Equation~\eqref{eq:auroc-policy-gamma} is the ordinary paired
AUROC replacement assessment under the declared replication challenge. As
$\Gamma$ grows, it asks how far that same conclusion survives increasingly
concentrated recombinations of the cases represented in each evaluation.

\subsection{The empirical case-mix breakdown factor}

The declared procedure locates the first value of $\Gamma$ at which either
the supported-magnitude requirement or the literal-survival requirement fails.
The reported value is therefore fixed by the search rule.

\begin{definition}[Empirical case-mix breakdown factor]
For the oriented comparison of candidate $A$ with incumbent $B$, define
\begin{equation}
\Gamma_{\mathrm{ECM},A:B}^{\dagger}(\mathbf D)
=
\inf\left\{
\Gamma\geq1:
\theta_{A:B;K}(\Gamma;\mathbf D)\leq\delta
\ \text{or}\ 
\psi_{K,d}^{A:B}(\Gamma;\mathbf D)\leq\gamma
\right\}.
\label{eq:auroc-ecm-breakdown}
\end{equation}
\end{definition}

The quantity is the smallest common upper within-class concentration factor at
which the replacement justification first fails. Equivalently, before the
breakdown, the conclusion survives every pair of empirical positive- and
negative-class redistributions whose individual weights are capped by
$\Gamma/N_+$ and $\Gamma/N_-$. A larger value means that a stronger empirical
concentration challenge is required to defeat the declared replacement policy.

If Equation~\eqref{eq:auroc-policy-gamma} already fails at $\Gamma=1$, then
$\Gamma_{\mathrm{ECM},A:B}^{\dagger}(\mathbf D)=1$: there is no supported replacement
conclusion before case-mix stress is applied. If the failure set in
Equation~\eqref{eq:auroc-ecm-breakdown} is empty, its infimum is $+\infty$ and
the report should state that no breakdown occurred over the full range of
redistributions permitted by the observed cases. In a finite evaluation the weight sets eventually saturate, because a
cap of at least one permits unrestricted simplex weights. A zero crossing of
$\theta_{A:B;K}(\Gamma;\mathbf D)$ may be informative as a secondary sensitivity
diagnostic, but the policy breakdown is primary because it preserves the
paper's practical-magnitude and literal-survival requirements.

The breakdown factor remains indexed by the declared assessment: the paired
models, observed evaluations and case sets, evaluation regime, replication
order, and policy thresholds $\delta$, $d$, and $\gamma$. Its reproducibility
requires the weight family and search rule to be fixed before the result is seen. The quantity
remains dependent on the declared context.

\subsection{Computation and breakdown reporting}

Four counts have different roles. Let $N_m$ be the number of evaluation units
in evaluation $D_m$, let $M$ be the number of separately conducted admissible
empirical evaluations, let $J$ be the number of computational empirical-resampling
replications used when only one evaluation is observed, and let $K$ be the
declared order of the replication challenge. Increasing $M$ adds
empirical evidence, increasing $J$ reduces Monte Carlo error in a projected
analysis, and increasing $K$ changes the criterion by placing greater emphasis
on adverse evaluations. Observed support requires $M\geq K$, while projected
support requires $J\geq K$; otherwise the corresponding order-$K$ assessment
is unavailable.

\paragraph{Several empirical evaluations.}
Suppose $D_1,\ldots,D_M$ are separately conducted full evaluations satisfying
$\Regime$. At each searched concentration factor, compute
\[
z_m(\Gamma)
=
Z_\Gamma^{A:B}(D_m),
\qquad m=1,\ldots,M.
\]
Applying the supported-magnitude and literal-survival summaries in
Equations~\eqref{eq:observed-empirical-support} and
\eqref{eq:literal-survival-estimator} directly to this indexed effect list
gives
\[
\theta_{A:B;K,M}^{\,\mathrm{obs}}(\Gamma)
\qquad\text{and}\qquad
\psi_{K,M,d}^{A:B,\mathrm{obs}}(\Gamma).
\]
The observed-replication breakdown factor is
\begin{equation}
\Gamma_{\mathrm{ECM},A:B}^{\dagger,\mathrm{obs}}
=
\inf\left\{
\Gamma\geq1:
\theta_{A:B;K,M}^{\,\mathrm{obs}}(\Gamma)\leq\delta
\ \text{or}\ 
\psi_{K,M,d}^{A:B,\mathrm{obs}}(\Gamma)\leq\gamma
\right\}.
\label{eq:auroc-ecm-breakdown-observed}
\end{equation}

\paragraph{One observed evaluation.}
When only one observed full evaluation $D_{\mathrm{obs}}$ is available,
generate $J$ computational replications $D^{(1)},\ldots,D^{(J)}$ under the
declared empirical-resampling procedure. Every computational replication must preserve the paired scores,
follow that procedure, and repeat the optimization in
Equation~\eqref{eq:auroc-z-gamma}. For every computational replication and
searched concentration factor, compute
\[
z_j(\Gamma)
=
Z_\Gamma^{A:B}\!\left(D^{(j)}\right),
\qquad j=1,\ldots,J.
\]
The same $D^{(j)}$ must be reused throughout the $\Gamma$ search. Changes in
$z_j(\Gamma)$ then arise from relaxing the concentration bound rather than from
using different computational resamples at different values of $\Gamma$.

Applying the two order-$K$ finite-list functionals to this indexed effect list
gives
\[
\widehat\theta_{A:B;K,J}^{\,\mathrm{proj}}(\Gamma)
\qquad\text{and}\qquad
\widehat\psi_{K,J,d}^{A:B,\mathrm{proj}}(\Gamma).
\]
The first value of $\Gamma$ at which either approximated profile crosses its policy
threshold gives the projected empirical case-mix breakdown approximation
\begin{equation}
\widehat\Gamma_{\mathrm{ECM},A:B}^{\dagger,\mathrm{proj}}
=
\inf\left\{
\Gamma\geq1:
\widehat\theta_{A:B;K,J}^{\,\mathrm{proj}}(\Gamma)\leq\delta
\ \text{or}\ 
\widehat\psi_{K,J,d}^{A:B,\mathrm{proj}}(\Gamma)\leq\gamma
\right\}.
\label{eq:auroc-ecm-breakdown-projected}
\end{equation}
The $J$ computational replications describe the evaluation effects produced by
the declared resampling procedure. They do not supply $M=J$ observed
evaluations. Increasing $J$ reduces Monte Carlo error but cannot expose a case
type, environment, measurement process, or failure mechanism absent from
$D_{\mathrm{obs}}$. The resulting profiles and breakdown factor must therefore
be described as projected.

For the observed path, the first failure point is a finite-list functional,
subject only to numerical optimization error. For the projected path, finite
$J$ adds Monte Carlo approximation error under the constructed resampling law.
Neither breakdown factor is an inferential guarantee about unseen cases or
contexts. The projected calculation should be repeated across random seeds and
larger values of $J$ to assess stability of the first failure point.



\subsection{Numerical verification and minimum reporting}

Writing $G=(G_{ik})$, Equation~\eqref{eq:auroc-z-gamma} minimizes the bilinear
form $w_+^\top G w_-$ over two capped simplices. A feasible weighting supplies
an upper bound on the infimum and can demonstrate failure when its value
violates a policy requirement. Numerically verifying that the requirement
continues to hold needs a valid global lower bound; alternating optimization
without such a bound may stop above the true infimum and overstate robustness.

Monotonicity permits a bracketed search for the corresponding observed or
projected breakdown factor rather than evaluation of a dense curve solely for
reporting. The numerical tolerance, global-optimization gap, Monte Carlo error,
and seed and $J$ stability should be recorded. The full profiles remain
useful diagnostics and should be retained by the software, but they need not
appear in the headline report.

A minimum report gives the decision results at $\Gamma=1$, the computed
breakdown factor, its numerical tolerance and optimization gap, Monte Carlo
error where applicable, and which of the supported-magnitude and
literal-survival requirements fails first. Every breakdown factor must be
identified as observed-replication or projected. After the evidence source has
been made explicit, unsuffixed notation may be used for readability. A possible
statement is:
\begin{quote}
At the observed case mix, model $A$ satisfied the order-$K$ supported-magnitude
and literal-survival requirements for replacement of model $B$. Under the
symmetric capped empirical case-mix challenge, the estimated breakdown
factor was $\widehat\Gamma_{\mathrm{ECM},A:B}^{\dagger}=[\text{value}]$,
computed with numerical tolerance [value] and global-optimization gap [value].
Through that factor, no observation could receive more than [value] times its
uniform within-class weight, and every admissible weighting retained effective
support of at least [fraction] of each observed outcome class. The first
failure involved [supported magnitude, literal survival, or both]. This is a
finite-evidence robustness result over represented cases, not a confidence
statement about unseen cases or contexts.
\end{quote}
With several empirical evaluations, the statement should open, ``Across $M$
separately conducted empirical evaluations satisfying the declared regime.''
With one resampled evaluation, it should open, ``Under the declared fixed-model
empirical-resampling procedure,'' and the breakdown factor in the statement
must be qualified as projected; Monte Carlo error and stability across seeds
and values of $J$ must also be reported.

\subsection{Scope of the empirical challenge}

The empirical target consists of redistributions of the case types contained
in the evaluation. Every admissible mix recombines observed cases. A type absent
from the evaluation receives weight zero at every $\Gamma$; reaching it requires
new data or an explicit extrapolation model. Individual reweighting may also
concentrate on sample-specific errors instead of a named scientific
subpopulation. The breakdown factor therefore has no interpretation as the
magnitude of a known deployment shift.

The AP construction has the same boundary. Prior standardization changes the
class prior while retaining the observed class-conditional score behavior.
Computational replication resamples the observed units. Empirical case-mix
reweighting redistributes those units within outcome class. These procedures
cannot reveal score behavior or failure mechanisms absent from the observed
evaluations. Separate empirical evaluations can expose such variation when it
falls within the declared regime.

No value encountered in the $\Gamma$ search is interpreted as an external
target population. The procedure starts from uniform empirical weights,
increases a normalized concentration allowance according to a fixed symmetric
rule, and reports where the declared replacement policy breaks. The empirical
case-mix breakdown factor is a reproducible property of the paired comparison
under the declared assessment. It measures robustness over represented cases
and remains silent about absent case types and mechanisms.


\begin{thebibliography}{99}
\setlength{\itemsep}{0.25em}

\bibitem{acerbi2002spectral}
Acerbi, C. (2002).
Spectral measures of risk: A coherent representation of subjective risk aversion.
\emph{Journal of Banking \& Finance}, 26(7), 1505--1518.

\bibitem{bareinboim2013transportability}
Bareinboim, E., \& Pearl, J. (2013).
A general algorithm for deciding transportability of experimental results.
\emph{Journal of Causal Inference}, 1(1), 107--134.

\bibitem{bestgen2015random}
Bestgen, Y. (2015).
Exact expected average precision of the random baseline for system evaluation.
\emph{The Prague Bulletin of Mathematical Linguistics}, 103, 131--138.

\bibitem{bental2013robust}
Ben-Tal, A., den Hertog, D., De Waegenaere, A., Melenberg, B., \& Rennen, G. (2013).
Robust solutions of optimization problems affected by uncertain probabilities.
\emph{Management Science}, 59(2), 341--357.

\bibitem{chernozhukov2013intersection}
Chernozhukov, V., Lee, S., \& Rosen, A. M. (2013).
Intersection bounds: Estimation and inference.
\emph{Econometrica}, 81(2), 667--737.

\bibitem{cohen1960coefficient}
Cohen, J. (1960).
A coefficient of agreement for nominal scales.
\emph{Educational and Psychological Measurement}, 20, 37--46.

\bibitem{elkan2001foundations}
Elkan, C. (2001).
The foundations of cost-sensitive learning.
In \emph{Proceedings of the Seventeenth International Joint Conference on Artificial Intelligence}, 973--978.

\bibitem{duchi2021statistics}
Duchi, J. C., Glynn, P. W., \& Namkoong, H. (2021).
Statistics of robust optimization: A generalized empirical likelihood approach.
\emph{Mathematics of Operations Research}, 46(3), 946--969.

\bibitem{lipton2018detecting}
Lipton, Z. C., Wang, Y.-X., \& Smola, A. (2018).
Detecting and correcting for label shift with black box predictors.
In \emph{Proceedings of the 35th International Conference on Machine Learning}, PMLR 80, 3122--3130.

\bibitem{popper1959logic}
Popper, K. R. (1959).
\emph{The Logic of Scientific Discovery}.
Hutchinson.

\bibitem{popper1963conjectures}
Popper, K. R. (1963).
\emph{Conjectures and Refutations: The Growth of Scientific Knowledge}.
Routledge.

\bibitem{saerens2002adjusting}
Saerens, M., Latinne, M., \& Decaestecker, C. (2002).
Adjusting the outputs of a classifier to new a priori probabilities.
\emph{Neural Computation}, 14(1), 21--41.

\bibitem{siblini2020calibration}
Siblini, W., Fr\'ery, J., He-Guelton, L., Obl\'e, F., \& Wang, Y. (2020).
Master your metrics with calibration.
In \emph{Advances in Intelligent Data Analysis XVIII}, 457--469.

\bibitem{shapiro2017distributionally}
Shapiro, A. (2017).
Distributionally robust stochastic programming.
\emph{SIAM Journal on Optimization}, 27(4), 2258--2275.

\bibitem{storkey2009training}
Storkey, A. (2009).
When training and test sets are different: characterizing learning transfer.
In \emph{Dataset Shift in Machine Learning}, 3--28. MIT Press.

\bibitem{wang1996premium}
Wang, S. (1996).
Premium calculation by transforming the layer premium density.
\emph{ASTIN Bulletin}, 26(1), 71--92.

\end{thebibliography}

\end{document}
